\chapter{Introduction}\label{ch:introduction}

\section{Motivation}\label{sec:motivation}

Large language models increasingly run inside software pipelines rather
than in front of a person. A task is broken into sub-tasks, each
sub-task is dispatched somewhere, and the dispatch target can be
another model --- HuggingGPT uses one model as a controller that routes
each sub-task to a specialised one~\parencite{shen2023hugginggpt}.
Every such decision needs to know which model is best at this kind of
work, and leaderboards do not say. They
report one aggregate score per model, averaged over tasks that have
little to do with one another~\parencite{boubdir2024elo}, and two
models with the same average can be strong in different
places~\parencite{raju2024constructing}. A router holding only the
average sends both the same way.

The published breakdowns do not go much deeper. MMLU-Pro, the corpus
this thesis uses, groups its questions into fourteen editorial
categories~\parencite{wang2024mmlupro}, inherited from
MMLU~\parencite{hendrycks2021measuring}, and no finer per-model
breakdown of it is published. ``Mathematics'' is one column covering
combinatorics, real analysis and elementary probability alike. If two
models diverge \emph{inside} mathematics, nothing on any leaderboard
records it. That is the gap this thesis works in. What it builds is
the measurement such a map would rest on --- not the router that
would consume it; whether the measurement can be trusted is the
question the thesis actually answers.

Getting a finer reading means judging model outputs at a scale where
expert human annotation is not
affordable~\parencite{zheng2023judging}. The tractable substitute is an
LLM judge, and the tractable form of the judgment is pairwise. Asking
which of two answers is better for one specific query avoids the
calibration problem that scalar scoring runs into, since a judge has no
stable internal reference for what a score of 7 means across domains,
question types and model generations~\parencite{zheng2023judging}.
Chatbot Arena runs that format at scale, with human voters,
Bradley--Terry aggregation and a public
leaderboard~\parencite{chiang2024chatbotarena}. Its prompts are
free-form and carry no answer key, so there is nothing to check a
vote against. This thesis judges a corpus
that has a key. Every verdict can therefore be checked, and that turns
out to decide the result.

\subsection*{The argument, stated in advance}

A judge should do better when it is told what ``better'' means
\emph{here}. A rubric that has to cover all of computer science can
only say domain-general things --- be correct, be clear, show your work
--- because every criterion sharp enough to discriminate is false of
part of the field. A rubric that covers only time-series model
diagnostics can name which stationarity checks are load-bearing and
which residual diagnostics decide the case. TaxoArena therefore splits
an evaluation corpus into small, topically coherent cells, writes a
rubric for each cell, and has an LLM judge compare model answers
pairwise inside cells. A narrow cell makes a specific rubric writable,
and a specific rubric was supposed to make a better judge.

The checks show the judge does not decide that way. It works out which
answer is correct and picks that side. Across the settled eight-domain
batch, replacing the cell rubric with a generic prompt moves only how
often the judge abstains
(Section~\ref{sec:res-c5}); questions routed off-subject are judged at
the same agreement rate as on-subject ones
(Section~\ref{sec:res-offsubject}); and the tie rate nearly doubles
exactly where the answer key stops discriminating
(Section~\ref{sec:res-tie-signature}). The rubric changes almost no
outcome, because on a keyed corpus a capable judge has a shortcut and
takes it.

What remains of the partition's claim is narrower than the thesis set
out to test. A pre-registered analysis found that cells do not carry
different true model rankings on this corpus and roster
(Section~\ref{sec:res-granularity}), so a per-cell ranking has little
to add to a global one. The paired arena runs of
Chapter~\ref{ch:results} therefore answer a smaller question: whether
the cell-scoped rubric changes what judging concludes at all. It does
not. The two arms rank identically in seven of eight domains and
within one adjacent swap in the eighth; what the rubric changes is how
often the judge abstains, and that difference belongs to the judging
condition as a bundle rather than to the rubric text alone
(Section~\ref{sec:res-c5}).

The result is a boundary. Wherever correctness is checkable, a capable
judge will shortcut the rubric, so on this corpus a cell-scoped rubric
had little left to improve. The setting where it could matter ---
open-ended tasks with no verifiable key --- is exactly the routing
setting, and it is the one this corpus cannot test.

\subsection*{What a partition is for}\label{sec:problem-statement}

\textbf{A cell is a scope for a rubric.} A semantic partition of an
evaluation corpus is usually a reporting device, deciding which bucket
each score is written into. Here it does something else, and the
difference decides what this thesis can claim: splitting a cell
reduces the surface over which a judge must generalise, which is the
whole mechanism, and no amount of prompt enlargement substitutes for
it --- narrowing the cell is what changes what the criteria's
intersection contains.

\textbf{Homogeneity is supposed to buy three things.} First, a rubric
becomes writable in the sense above. Second, pairwise outcomes within
the cell come closer to exchangeable, because the queries a comparison
could have been drawn from resemble one another; the Bradley--Terry
model that aggregates the verdicts assumes a single latent strength per
model within the cell, and that assumption is a statement about the
cell's homogeneity rather than about the models. Third, the per-cell
ranking is allowed to differ from the global one, which is the only way
a capability \emph{profile} carries more information than a capability
\emph{score}.

The third is the one that turned out to be empirical rather than
architectural, and it is the one the data refused
(Section~\ref{sec:res-granularity}). This section is an argument, not a
measurement. Chapter~\ref{ch:results} reports what happened when each
part of it was tested.

\section{The System in Outline}\label{sec:approach}

TaxoArena works in three stages. The fourteen MMLU-Pro domain labels
seed top-level anchor nodes, which are then refined by embedding
geometry: queries migrate to geometrically closer anchors, and each
anchor is recursively split into a tree of coherent leaves. Each leaf
then induces a judge rubric from construction-set queries and their
correct answers. Finally, held-out model responses are compared
pairwise within leaves and aggregated by a Bradley--Terry
model~\parencite{bradley1952rank,hunter2004mm} into per-cell rankings.

\textbf{Position of the taxonomy.} The induced taxonomy is \emph{not}
the contribution of this thesis. It is the instrument by which
per-domain ranking resolution is made available. The main text
therefore presents the construction only at the level needed to
establish that its cells are granular, coherent and stable enough for
the arena to mean anything (Section~\ref{sec:arch-requirements}); the
formalism, the parameter derivations and the fixed-point certificate
are in Appendix~\ref{app:dag-formalism} and
Appendix~\ref{app:numerics}. That ordering is forced by a measurement
rather than chosen late. An elevenfold change in the number of cells,
from 14 to 152, moves the system's measured discriminative power by
less than the analysis can resolve
(Section~\ref{sec:res-granularity}). A construction whose tuning has no
measured evaluative consequence is an instrument, and it is documented
like one.

\textbf{Reference-informed at construction, answer-key-blind at test.}
The judge-generation phase uses correct MMLU-Pro answers from the
construction query set to induce the rubrics, so the system is
\emph{reference-informed} at construction time. The arena phase judges
held-out responses without access to those queries' correct answers, so
the evaluation is \emph{answer-key-blind} at test time; the MMLU-Pro
accuracy oracle is revealed only after all matches are complete. A
hold-out filter enforces the separation in code, and a build-time
assertion checks the filter against the membership test itself
(Section~\ref{sec:arch-judge}).

Answer-key-blind is not answer-blind, and the difference is
load-bearing for Chapter~\ref{ch:results}; the mechanism, and what
the judge is shown mode by mode, is Section~\ref{sec:arch-judge}.

\textbf{The title terms.} \emph{Dynamic Domains} refers to boundaries
seeded from MMLU-Pro labels and then adapted by geometric evidence, so
the evaluation dimensions follow the data rather than a fixed editorial
list. \emph{Query-Adaptive} names the corpus-level sense of that
adaptivity: the domain \emph{structure} adapts to the query corpus as a
whole, which is distinct from the per-prompt sense in which
Prompt-to-Leaderboard conditions a full ranking on each individual
prompt~\parencite{frick2025p2l}. \emph{Multidimensional} means the
output is one ranking per cell rather than one scalar per model.

\section{Three Links and the Research Questions}\label{sec:research-questions}

The claim that a partition improves evaluation is not one claim. It is
a chain of three claims, and each can fail on its own:

\begin{center}
coherent cells
$\;\xrightarrow{\;\text{(1)}\;}\;$ specific rubrics
$\;\xrightarrow{\;\text{(2)}\;}\;$ better judgments
$\;\xrightarrow{\;\text{(3)}\;}\;$ ranking closer to ground truth.
\end{center}

Link~1 asks whether the construction produces cells homogeneous enough
to write a rubric about. Link~2 asks whether cell-scoped rubrics are in
fact more specific than generic ones. Link~3 asks whether a judge
holding a specific rubric produces verdicts that track the truth more
closely. Naming the links separately is what makes this thesis's
results readable, because its negative results land on different links
and mean different things there.

The decomposition was written after the pre-registered granularity
analysis came back against the premise of the original second research
question --- that cells carry different true model rankings. A
contrast built on that premise measures a difference the data gives no
room to exist, so the question was replaced rather than repaired, and
the replacement is stated here so the rewrite is on the record rather
than left to be discovered. What the null removes is one specific
thing the partition could have bought. It leaves the other two links
standing.

\begin{enumerate}

    \item[\textbf{RQ1.}] \emph{What decides the verdicts} of
    answer-key-blind, domain-conditioned pairwise LLM judgments --- and
    do those judgments recover per-domain capability rankings on
    held-out queries?
    \textit{Primary metric: agreement between the judge's winner and
    the MMLU-Pro answer key, per verdict. Spearman $\rho$ against
    ground-truth accuracy is reported alongside, under both tie
    policies, with its quantisation stated.}

    \item[\textbf{RQ2.}] Does a cell-scoped rubric enable \emph{better
    judging} than judging the same questions under a generic rubric?
    \textit{Primary metric: per-verdict agreement between the judge's
    winner and the answer key, compared between the two arms on an
    identical question set. Spearman $\rho$ is reported alongside, as
    $\Delta\rho = \rho_{\textsc{main}} - \rho_{\textsc{c5}}$, so a
    positive value favours cell-scoped judging --- reported for
    continuity, not relied on: the registered decisive threshold is
    withdrawn (Section~\ref{sec:exp-metrics}). RQ2 varies the rubric,
    not the partition: both arms judge the same questions in the same
    cells under the replay design of
    Chapter~\ref{ch:experimental-design}.}

\end{enumerate}

RQ2 is answerable with the same arena, and it is the question the
system actually poses. The change from the original wording matters for
what gets reported: agreement between judge and ground truth
\emph{within} a grouping, rather than rank divergence \emph{between}
groupings.

\section{Contributions}\label{sec:research-goal}

\begin{enumerate}

    \item \textbf{A measured account of what an LLM judge decides on a
    corpus with a verifiable answer key.} This
    answers RQ1: the verdicts track answer correctness rather than
    reasoning quality, recovering the key's ordering on $196$ of the
    $200$ decidable model pairs across eight domains. Each supporting
    measurement carries a control sharing every confound with its
    treatment except the one varied (Section~\ref{sec:res-judge}).

    \item \textbf{A pre-registered rubric-specificity baseline.}
    This answers RQ2 and tests link~2. Paired arms
    judged identical comparisons under cell-scoped and generic rubrics,
    with reading rules fixed in advance. The rubric does not change
    what judging concludes --- the arms rank identically up to one
    adjacent swap --- and the abstention difference that remains
    belongs to the judging condition as a bundle
    (Section~\ref{sec:res-c5}).

    \item \textbf{A construction method as an instrument, with its
    limitations.} Link~1 rests on it: the induced
    87 leaves score $J = 0.2538$ against $0.1270$ for the editorial
    categories and $0.1260$ for a random 88-cell sub-split, so the
    gain is not bought by cell count; the construction reaches a
    certified fixed point (Section~\ref{sec:arch-requirements}). It is
    seed-dependent --- median ARI $0.677$ across twenty sweeps --- so a
    query's leaf is not a seed-invariant fact.

    \item \textbf{A pre-registered null on whether finer cells rank
    models differently at all.} Link~3's premise,
    and the premise of a capability \emph{profile}: between-cell rank
    correlation barely moves from 14 to 152 cells, so the answer is a
    null, readable because its outcomes were named in advance. Power
    and roster bound its reach (Section~\ref{sec:res-granularity}).

    \item \textbf{The three-link decomposition, offered as a way of
    stating the claim.} Splitting ``a taxonomy
    improves evaluation'' into three links is what lets this thesis
    say which link each negative result lands on. A framing, not a
    finding, and tagged accordingly.

\end{enumerate}

The thesis set out to test one claim: that a judge built from
guidelines extracted from a homogeneous set of queries tells you more
about what a model can do than a general-purpose judge does. Answering
that meant building a partition first, and most of the engineering went
there. The claim was always about the judge, which is why
Chapter~\ref{ch:results} leads with what the judge does and only then
with what the partition does. An active matchmaking scheduler
(Section~\ref{sec:meth-scheduling}) keeps judge-call cost tractable; it
supports the arena rather than standing as a contribution, and its
behaviour changed mid-project in a way that bounds what the paired runs
can be read to say.

\section{Scope and Validity}\label{sec:intro-scope}

Six conditions bound everything downstream, and they are stated before
any result is read.

\textbf{One --- the corpus.} MMLU-Pro is multiple-choice with a
verifiable key. That is what makes comparison against ground truth
possible, and it is the same property that lets a capable judge bypass
the rubric by verifying the answer. Both are true.
Chapter~\ref{ch:discussion} treats the second as a measurement of where
the architecture binds rather than as a defect.

\textbf{Two --- selected domains.} The arena is a selected-domain
study, not a global validation: judge-call budget was spent on eight
of the fourteen MMLU-Pro domains at high statistical power rather
than spread across all of them (selection rule in
Chapter~\ref{ch:experimental-design}). Refit at the settled roster,
the screen behind that selection saturates, so the eight domains are
read as replications of one measurement rather than as screened
contrasts (Section~\ref{sec:res-law}).

\textbf{Three --- quantisation.} At eight models the Spearman
statistic moves in steps of $0.0119$, so absolute rank correlations
are not stable to the third decimal and this thesis does not quote
them that way (Section~\ref{sec:res-tie-policy}).

\textbf{Four --- the scheduler seed.} A seed replicate moved the
paired $\Delta\rho$ contrast by six grid steps against a registered
decisive threshold of two, since withdrawn
(Section~\ref{sec:exp-metrics}). A rank-correlation magnitude is a
property of the schedule as much as of the arms, and this thesis does
not read one as the latter (Section~\ref{sec:res-tie-policy}).

\textbf{Five --- no intervals.} No Bradley--Terry confidence interval
appears anywhere in this thesis. The variance computation was wrong
before a dated fix and the intervals were never recomputed, so there is
nothing to print (Section~\ref{sec:disc-debt}).

\textbf{Six --- two earlier run generations are withdrawn.} A pilot
arena did not persist response text for four of its eight models, so
half of its comparisons were decided by which side had text at all;
and a twelve-model paired campaign ran before two scheduler repairs
that the pilot's defect analysis forced. Neither contributes a
number to this thesis. The defect, the repair, and what the
withdrawals cost are reported as results
(Section~\ref{sec:res-capture-gap}).

The evaluation's validity rests on the reference-informed,
answer-key-blind separation defined in Section~\ref{sec:approach}. The
70/30 construction/arena-test split and the information-separation
protocol that enforces it are specified in
Chapter~\ref{ch:experimental-design}.

\section{Thesis Structure}\label{sec:thesis-structure}

Chapter~\ref{ch:literature-review} gives the background the argument
uses: what is known to go wrong with LLM judges, pairwise ranking and
Bradley--Terry estimation, the family the construction belongs to, and
query routing as the consumer of the output.
Chapter~\ref{ch:methodology} is the method. It opens by stating what a
partition must satisfy for the arena to mean anything and shows that
the induced partition satisfies it, then specifies the judge, the
estimator and the scheduler, and closes with the measurement rules
the thesis holds itself to.
Chapter~\ref{ch:experimental-design} fixes the corpus, the rosters, the
conditions, the pre-registrations with their dates, and the statistical
protocol.
Chapter~\ref{ch:results} reports what was measured, in mechanism order:
first what the judge does, then what the partition does.
Chapter~\ref{ch:discussion} interprets those results, states where the
architecture binds, records the errors the apparatus caught in its own
published numbers, and concludes.
